
\graphicspath{{Figures/Discussion}}

\chapter{Discussion}\label{ch:discc}
\begin{itemize}
  \item \textbf{The Reminder:} Re-state the core biological/computational problem your thesis tackled
  \item \textbf{The Headline Finding:} State your absolute biggest achievement without listing a wall of specific metrics.
  \item \textbf{The Executive "Why":} Give a one-sentence conceptual reason for your success.
\end{itemize}

\section{Model Performance}
\begin{itemize}
  \item Why are they this way? Explain the mechanics behind the scores
  \item Where does it fit? Compare this to the literature.
\end{itemize}

\section{Ablation Study and Biological Significance}
\begin{itemize}
  \item Why are they this way? Explain the mechanics behind the scores
  \item Where does it fit? Compare this to the literature.
\end{itemize}

\section{Candidate Variants}

The $katG$ C>G SNP has biological grounding, as $katG$ encodes a catalase-peroxidase that activates the prodrug for isoniazid [CITE]. With non-synonymous variants in $katG$ being a substantial reservoir for resistance-conferring variants in clinical samples, an uncatalogued variant that is both significant and associated with resistance, and produced a $-6.26\%$ change in sensitivity when ablated poses as a high-priority variant warranting further characterisation and investigation.  

The `gene-not-found' C>T variant also showed a considerable drop in sensitivity of $4.56\%$, and has potential for investigation as a resistance-marker. This variant is located 14bp upstream of $fabG1$, an enzyme involved in mycolic acid biosynthesis, and more importantly, a gene that contains resistance-conferring variants for isoniazid and ethionamide \cite{valafar2021systematic}. Given an approximate 5$\%$ sensitivity dip post-ablation, and proximity to $fabG1$, this variant warrants further investigation. 

Lastly, an $rpoB$ C>T SNP yielded a drop in sensitivity of 0.73$\%$. This is very likely a confounder, as $rpoB$ encodes the $\beta$ subunit of RNA polymerase [CITE], which is the target for rifampicin. This appearance is likely because MDR-TB strains carry resistance to both rifampicin and isoniazid. Future weight matrix revisions for isoniazid warrants removal to prevent FPs. 

This validates our internal pipelines, which identified the correct gene and location. The enrichment of $rpoB$ C>T annotated SNP, is one of the strongest candidates for validation from the FN subset. All $rpoB$ variants close to, or within RRDR (codons 426-452). 


Multiple SNPs at distinct positions, indicates possible diverse resistant variants over a single dominant mutation. This is consistent with literature surrounding clincinal ethambutol resistance which report a wide spectrum of $embB$ variants. These candidates are of high-priority for validation. Furthermore, variants appeared in auxiliary genes involved in $embB$ resistance such as $embC$ and $ubiA$.
\section{Limitations}
\begin{itemize}
  \item Why are they this way? Explain the mechanics behind the scores
  \item Where does it fit? Compare this to the literature.
\end{itemize}
\section{Future Studies}
\begin{itemize}
  \item Why are they this way? Explain the mechanics behind the scores
  \item Where does it fit? Compare this to the literature.
\end{itemize}
